\documentclass[11pt]{letter}
\usepackage[utf8]{inputenc}
\usepackage[T1]{fontenc}
\usepackage[margin=1in]{geometry}
\usepackage{hyperref}
\usepackage{siunitx}
\urlstyle{same}

\signature{Jonathan Mace\\Corresponding author}
\address{Jonathan Mace\\Microsoft Research\\Redmond, WA\\jonathanmace@microsoft.com}
\date{\today}

\begin{document}

\begin{letter}{%
Editor-in-Chief\\
\textit{Journal of the Acoustical Society of America}}

\opening{Dear Editor-in-Chief,}
\small

Please consider our manuscript, ``What pitch is a growling stomach? A unified multi-mode acoustic model of borborygmi,'' for publication in \textit{Journal of the Acoustical Society of America}. The paper develops a first-principles analytical framework for bowel sounds by combining four candidate resonance mechanisms: tissue-constrained bubble resonance, Helmholtz resonance through constrictions, axial piston modes, and radial breathing modes of gas columns in elastic tubes. Using measurable anatomical and material parameters rather than empirical fitting, the model predicts that the tissue-constrained bubble mechanism occupies much of the clinically observed healthy bowel-sound band, while the other mechanisms dominate in more limited regions of parameter space.

We believe the manuscript is a strong fit for \textit{JASA} because it connects acoustic theory directly to a familiar but poorly understood biomedical sound. The contribution is not a classifier or phenomenological survey; it is an acoustical model that explains why particular frequency bands arise and how spectral changes can follow from gas-pocket volume, lumen geometry, or constrictions. In that sense, the paper should interest readers in biomedical acoustics, bubble dynamics, physiological sound analysis, and non-invasive diagnostic monitoring. It also forms part of our broader research programme on acoustics and vibroacoustics in fluid-filled and gas-containing biological systems.

The manuscript is original, has not been published previously, is not under consideration elsewhere, and has been approved by both authors. Brian R. Mace's submission email is \texttt{b.mace@auckland.ac.nz}. We hope the paper will be of interest to \textit{JASA} as a quantitatively grounded biomedical-acoustics contribution with clear diagnostic relevance.

\closing{Yours sincerely,}

\end{letter}
\end{document}
